\subsection{The actual coefficient map under weight Frobenius saturation}

We opened Bhatt--Lurie--Mathew's author PDF
\cite{Root_BhattLurieMathew2020}: Construction~2.1.3 and
Remark~2.1.4, page~13; Proposition~2.3.1, page~15;
Construction~2.5.1, page~18; and Definition~3.1.2, page~29.
For a $p$-torsionfree cochain complex its operator is
\[
 (\eta_pM)^q=\{x\in p^qM^q:dx\in p^{q+1}M^{q+1}\}.
 \tag{WF.1}
\]
A graded endomorphism $F$ with $dF=pFd$ determines
$\alpha_F:M\to\eta_pM$, $x\mapsto p^qF(x)$ in degree $q$.
Their saturation is the filtered colimit of the iterates
of this actual map. Their completion uses the quotients
by $V^rM+dV^rM$, where $FV=p$.
These are the precise constructions used in the following
calculation; no comparison with a differently specified
Frobenius is asserted.

\begin{proposition}[A-WF01: Proved; the weight Frobenius changes the coefficient comparison]
\label{AWF01}
Fix a prime $p$, put $R=\mathbb Z_{(p)}$ and
$S_R=\Gamma_R(T)$, with $T$ of weight one and form degree zero.
Let $A$ be a termwise torsionfree strict PD theory with
$q\leq w$, weight-zero algebra $R$ in every simplex,
and actual coefficient identification
$H^0(A_n^{w,*})=S_R^w$ in every weight and simplex.
This includes the standard admissible theory $G_R$.
Define, on homogeneous elements, the graded endomorphism
\[
                     F_w^q(a)=p^{w-q}a.             \tag{WF.2}
\]
It equips each simplex algebra with an actual Dieudonn\'e
algebra structure, functorially in the simplicial direction.
Its saturation and its natural coefficient map satisfy
\[
 H^0(\operatorname{Sat}_F(A_n))^w=
 \begin{cases}
 R,&w=0,\\
 \mathbb Q T^{[w]},&w>0,
 \end{cases}
 \qquad
 S_R^w\longrightarrow H^0(\operatorname{Sat}_F(A_n))^w
 \text{ is the indicated inclusion}.               \tag{WF.3}
\]
The natural composite to its completed saturation
\[
              A\longrightarrow\operatorname{Sat}_F(A)
                \longrightarrow\mathcal W\operatorname{Sat}_F(A)
 \tag{WF.4}
\]
sends every positive-weight coefficient $T^{[w]}$ to zero.
Consequently neither the saturation with its induced
coefficient map nor the completion (WF.4) is an admissible
replacement with that coefficient comparison.  This failure
already occurs for the actual admissible input $G_R$.
\end{proposition}

\begin{proof}
Support makes the exponent in (WF.2) nonnegative whenever
its source group is nonzero. Since weights and form degrees
add under products, (WF.2) is additive and multiplicative,
and sends the unit to itself. It commutes with every
simplicial map and with divided powers: the exponent on
$\gamma_j(a)$ is $j(w-q)$, exactly the exponent supplied
by the PD scalar rule on $\gamma_j(Fa)$.
The differential identity is the equality
\[
                 d(p^{w-q}a)
                   =p\,p^{w-q-1}da.                \tag{WF.5}
\]
When $q=w$, the group containing $da$ vanishes by support,
so that case also satisfies the identity without a
negative-power interpretation.

To verify the degree-zero Frobenius congruence, write a
zero-form as $a_0+a_+$, with $a_0\in R$ and $a_+$ in
positive weights. The latter belongs to the PD ideal, so
$a_+^p=p!\gamma_p(a_+)$ is divisible by $p$.
The binomial congruence and $R/pR=\mathbb F_p$ give
$(a_0+a_+)^p\equiv a_0\pmod p$, while (WF.2) gives the
same residue. This verifies all conditions of the cited
Dieudonn\'e-algebra definition. The associated actual
cochain map is
\[
                  \alpha_F|_{A^{w,q}}=p^w\operatorname{id}.
 \tag{WF.6}
\]
Its image belongs to (WF.1): $w\geq q$, and a nonzero
derivative has $w\geq q+1$.

For any nonnegative torsionfree cochain complex, the
closed degree-zero elements of $\eta_pM$ are exactly
those of $M$. They already satisfy every divisibility
condition in degree zero, and the ambient differential
is unchanged. Repeating this observation shows
\[
              H^0(\eta_p^rA_n)^w=S_R^w
                       \quad\hbox{for all }r\geq0. \tag{WF.7}
\]
There are no incoming degree-negative boundaries.
The maps on these closed groups in the saturation
sequence are multiplication by $p^w$, by (WF.6).
Filtered colimits of modules preserve kernels, so their
colimit computes the degree-zero cohomology of the
saturation. For $w=0$ this is the identity system on $R$.
For $w>0$ the explicit compatible identification of its
$r$th stage with the subgroup
$p^{-rw}R T^{[w]}\subset\mathbb Q T^{[w]}$ gives the
whole union $\mathbb Q T^{[w]}$.
It carries the initial coefficient map to inclusion.
This proves (WF.3) with its stated comparison map.

On the closed positive-weight subgroup of the saturation,
$F$ still acts by multiplication by $p^w$.
The defining uniqueness of $V$ and the equality $FV=p$
therefore give
\[
       V(c)=p^{1-w}c,
       \qquad
       V^r\bigl(p^{r(w-1)}c\bigr)=c
              \quad(c\in\mathbb Q T^{[w]}).         \tag{WF.8}
\]
Every argument of $V^r$ displayed here is a closed
coefficient element of the saturated complex.
Thus $c$ maps to zero in every quotient by
$\operatorname{im}(V^r)+\operatorname{im}(dV^r)$.
Its image in their inverse limit is zero, proving (WF.4)
on every positive-weight coefficient.

For the saturation alone, (WF.3) fails to be the required
coefficient isomorphism in every positive weight: $R\to\mathbb Q$
is not onto, as $1/p\notin R$.
For the completion, its induced coefficient map has a
nonzero element $T^{[w]}$ in its kernel for every $w>0$.
No eventual weight cutoff repairs either defect.
Finally the standard $G_R$ has the hypotheses by its
explicit integral PD identity-disk contractions,
nonnegative support and weight-zero description; the
same standard-model verification is rechecked in the
local-source construction of C09 used in this run.
This supplies the asserted actual admissible input.
\end{proof}
